\chapter{Introduction}

\begin{comment}

    Title
     - Fast launcher performance evaluation using reinforcement learning
     - Orbital ascent trajectory optimization using reinforcement learning
     - Orbital ascent trajectory optimization using reinforcement learning for the purpose of launcher performance evaluation
     - Fast launcher performance evaluation using reinforcement learning based guidance.
     - Fast orbital ascent trajectory optimization using reinforcement learning

    Describe problem(s) that we are trying to solve.
     - Launcher performance evaluation
     - Orbital ascent trajectory optimization
     - Slow performance of current launcher performance evaluation
     - Lifting body first stage
     - Unknown situations and circumstances

    
    Why it's relevant to the aerospace industry.
     - Space

    What the basic solution to this problem could be, current solutions
     - Semi empirical
     - Decomposition
     - Simplification by assumptions
     - Global optimization methods
    
    Research question

    Proposed solution of the preliminary thesis and why.
     - RL
     - Short description of what RL is capable of
     - Give examples of previous research or real world application (keep short or out)

    Short description on the approach that will be taken during the thesis to use reinforcement learning to solve the problem.

    Describe the contents in this report and how it will be used during the actual thesis.
     - Introduction into reinforcement learning
     - Tabular methods
     - Approximate solutions 
     - History of RL. Past, Present and Future (Or maybe just the latest in RL)
     - Thesis plan of attack

\end{comment}
As we delve deeper into our solar system it becomes increasingly difficult to navigate and control in the new uncharted worlds we are sending our spacecraft to. The non-uniform gravitational fields of asteroids, unexplored atmospheres, unexpected obstacles, unknown disturbances and unforeseen problems will all pose a challenge as deep space exploration missions become more ambitious.

If new, farther, more complex environments are to be explored in the future, spacecraft will need to become more autonomous and less reliant on ground personnel on Earth. Unknown problems requiring immediate solutions will need to be dealt with without human intervention because of the communication lags imposed by the physical limits of the speed of light. As an example this thesis investigates guidance and navigation of an orbital launch vehicle through an unfamiliar atmosphere.

A successful orbital insertion requires the launcher to carefully follow an optimized trajectory. This trajectory has a massive effect on the performance of the launcher, but unfortunately finding the optimal trajectory is not trivial. This problem falls under what is known as optimal control, a class of problems where the goal is to design a controller that optimizes a performance measure of a dynamical system's behavior through time, \cite{sutton2018reinforcement, betts2010practical}. 

A commonly used class of methods to compute orbital ascent trajectories is trajectory optimization. This is a set of mathematical techniques used to find the desired behavior of a dynamic system by finding the inputs as a function of time, \cite{betts2010practical, kelly2017introduction}. 

Using these methods to compute orbital ascent trajectories is computationally intensive and only yields open loop solutions for one particular trajectory. This is not an issue for normal launch operations on Earth, but it may not provide suitable solution when the flight conditions change, are different than unexpected or an unknown issue is encountered.

This thesis proposes the use of \acrfull{rl} as a method to compute closed-loop solutions for orbital ascent trajectory problems. With enough training an \acrshort{rl} agent is able to learn to solve complex and non-linear problems. Thus the goal is to develop an \acrshort{rl} agent that is able to compute orbital accent trajectories for a wide range of initial states, atmospheric conditions, launcher configurations and with the capability to generalize and successfully guide the launcher even through unknown environments.

Thus the main research objective and questions for this thesis are:

% \emph{“To develop a tool capable of computing large batches of orbital ascent trajectories for a wide range of initial conditions and launcher configurations, by taking advantage of Reinforcement Learning and investigate its effectiveness with respect to existing methods, using metrics such as the accuracy of the results and computational requirements.”}

% \emph{“To develop and train a Reinforcement Learning agent such that it provides a closed-loop orbital accent guidance law capable of generalizing beyond explored environments and compensate for unknown disturbances.”}

\emph{“To obtain a closed-loop orbital accent guidance law capable of generalizing beyond explored environments and compensate for unknown disturbances using Reinforcement Learning .”}


\textbf{How are orbital ascent trajectories currently computed?}
\begin{enumerate}
  \item What are the challenges faced when computing an orbital ascent trajectory?
  \item What methods are currently used or being researched to solve this problem?
  \item What are the limitations of methods?
\end{enumerate}

\textbf{What is \acrlong{rl}?}
\begin{enumerate}
  \item How does \acrshort{rl} work?
  \item What is \acrshort{rl} capable of?
  \item How has \acrshort{rl} been used so far?
  \item What state of the art technological developments make use of \acrshort{rl}?
\end{enumerate}

\todo { Add questions on using failed and successfully }
 
\textbf{How can \acrlong{rl} methods be used to solve the orbital ascent trajectory opti­mization problem?}
\begin{enumerate}
  \item What advantages do \acrshort{rl} methods have over existing methods?
  \item What properties of \acrshort{rl} can be used to circumvent the limitations of current orbital ascent trajectory optimization methods?
  \item Are \acrshort{rl} based orbital ascent trajectory computation methods capable of achieving the accuracy of existing methods?
  \item Is an \acrshort{rl} based closed-loop solution able to outperform a precomputed open-loop solution in case of unknown disturbances or errors in the training models.
\end{enumerate}

Chapter~2 starts with the theory of behind trajectory optimization and is directly followed by a review of the state of the art methods used in orbital ascent trajectory optimization.

Chapter~3 continuous with a history on \acrshort{rl} followed by a literature review of latest developments in \acrshort{rl}. This is then followed by several sections explaining the theory behind \acrshort{rl} ending with the more advanced methods that will be used throughout the main thesis project. The chapter then concludes with two examples demonstrating the different methods.

The next chapter contains a review of the state of the art research on \acrshort{rl} in spacecraft guidance, and finally the report is concluded with an overview of research questions and the proposed experiment to be conducted during the main thesis.

The use of \acrshort{rl} in spacecraft guidance, navigation and control opens up the possibilities to solve future problems as space exploration becomes more prevalent. For example thanks to its abilities to generalize and adapt in real time, it may be used in situations that are difficult to forsee or too cumbersome to manually take into account, such as hovering and manoeuvering around the non-uniform, rotating gravitational field of an asteroid, \cite{gaudet2012robust, willis2016reinforcement}, performing docking operations with non-cooperative target objects in cluttered environments, \cite{hovell2020deep, hovell2021deep, broida2019spacecraft}, or achieving precision planetary powered decent and landing on Earth, Mars and beyond \cite{gaudet2018deep}.
